\documentclass[11pt,a4paper]{article}
\usepackage[utf8]{inputenc}
\usepackage[T1]{fontenc}
\usepackage[french]{babel}
\usepackage{lmodern,geometry,microtype}
\usepackage{amsmath}
\usepackage{siunitx}
\usepackage[version=4]{mhchem}
\usepackage{xcolor,listings,hyperref}
\geometry{margin=2.4cm}
\sisetup{locale=FR,per-mode=symbol}
\DeclareSIUnit\molecule{molécule}
\DeclareSIUnit\atm{atm}
\lstset{language=Python,basicstyle=\ttfamily\small,backgroundcolor=\color{black!4},frame=single,breaklines=true,showstringspaces=false}
\hypersetup{colorlinks=true,linkcolor=blue!45!black,urlcolor=blue!55!black}

\title{Documentation de \texttt{1\_section\_efficace\_CO2\_15um.py}}
\author{Projet climat 2026}
\date{\today}

\begin{document}
\maketitle
\tableofcontents

\section{Rôle du programme}

Ce script est un prototype autonome consacré à la bande d'absorption du
\ce{CO2} autour de \SI{15}{\micro\metre}. Il télécharge les raies HITRAN,
calcule une section efficace avec un profil de Voigt, la convertit en unités
SI, puis l'affiche entre 12 et \SI{18}{\micro\metre}.

\section{Installation préalable}

Python~3.12 et les trois paquets de \path{requirements.txt} sont requis :

\begin{lstlisting}[language=bash]
python3.12 -m venv .venv
source .venv/bin/activate
python -m pip install -r requirements.txt
\end{lstlisting}

Les paquets externes sont \texttt{numpy}, \texttt{matplotlib} et
\texttt{hitran-api}. Attention : \texttt{hitran-api} s'importe sous le nom
\texttt{hapi}. \texttt{pathlib} appartient à la bibliothèque standard.
Le premier lancement exige une connexion au serveur HITRAN si la table locale
\texttt{CO2} n'est pas encore disponible.

\section{Imports du fichier}

\begin{lstlisting}
from pathlib import Path

import matplotlib.pyplot as plt
import numpy as np
from hapi import absorptionCoefficient_Voigt, db_begin, fetch
\end{lstlisting}

\begin{itemize}
  \item \texttt{Path} construit le chemin absolu du cache ;
  \item \texttt{numpy} crée les grilles et effectue l'interpolation ;
  \item \texttt{matplotlib.pyplot} produit le graphique ;
  \item \texttt{db\_begin}, \texttt{fetch} et
        \texttt{absorptionCoefficient\_Voigt} assurent l'accès à HITRAN et
        le calcul spectral.
\end{itemize}

\section{Constante globale}

\texttt{DOSSIER\_DONNEES} désigne \path{donnees_hitran/}, situé à côté du
script. HAPI y lit ou y crée \path{CO2.data} et \path{CO2.header}.

\section{Fonctions}

\subsection{\texttt{convertir\_lambda\_m\_en\_nombre\_onde\_cm\_1}}

\textbf{Entrée :} longueur d'onde scalaire ou tableau, en mètres.

\textbf{Retour :} nombre d'onde en \si{\per\centi\metre}, selon
\begin{equation}
  \widetilde\nu=\frac{1}{100\lambda}.
\end{equation}

\subsection{\texttt{charger\_section\_co2}}

Paramètres par défaut : 550--\SI{800}{\per\centi\metre}, pas de
\SI{0.05}{\per\centi\metre}, $T=\SI{296}{\kelvin}$ et
$p=\SI{1}{\atm}$. La fonction initialise la base, télécharge
l'isotopologue principal du \ce{CO2}, puis appelle le profil de Voigt.

\textbf{Retour :} \texttt{(nombres\_onde, sections\_m2)}, deux tableaux NumPy.
HAPI fournit initialement des \si{\centi\metre\squared\per\molecule}; le
facteur $10^{-4}$ donne des \si{\metre\squared\per\molecule}.

\subsection{\texttt{calculer\_section\_co2}}

La fonction convertit les longueurs d'onde demandées puis interpole la
section avec \texttt{numpy.interp}. Les valeurs hors de la plage HITRAN sont
fixées à zéro. Elle renvoie une section scalaire ou un tableau selon l'entrée.

\subsection{\texttt{tracer\_section\_co2}}

Cette fonction construit 3000 points de 12 à \SI{18}{\micro\metre}, calcule
les sections, puis ouvre un graphique. Elle constitue le point d'entrée lors
de l'exécution directe du fichier.

\section{Exécution et sortie}

\begin{lstlisting}[language=bash]
python 1_section_efficace_CO2_15um.py
\end{lstlisting}

La sortie est une fenêtre graphique ; aucun fichier n'est exporté. Le dossier
\path{donnees_hitran/} peut en revanche être créé ou complété. En environnement
sans affichage, il faut configurer un backend Matplotlib non interactif ou
remplacer \texttt{plt.show()} par une sauvegarde de figure.

\section{Place dans le modèle}

Ce prototype sert à valider l'installation de HAPI et les unités sur un seul
gaz. Le calcul multi-gaz utilisé par le modèle final se trouve dans
\path{2_sections_efficaces_gaz_hitran.py}.

\end{document}
